\section{Motociclos, ciclomotores e mini carros}

\subsection{Motociclos e ciclomotores até 2006}

Até 2006, os motociclos até 50~cm³ e os ciclomotores tinham uma matrícula
diferente da dos automóveis, atribuída pela câmara municipal. As placas tinham a
sequência \plate{0-XXX / 00-00}, em que \plate{XXX} é o código de três letras do
município onde era feito o registo (ver o anexo~\ref{sec:municipios}). Esta
regra foi fixada pelo Decreto Regulamentar n.º~13/98.

As chapas dos ciclomotores eram de fundo amarelo e as dos motociclos de
cilindrada inferior a 50~cm³ de fundo branco. Até à década de 1980, as chapas de
fundo amarelo eram também usadas por bicicletas e veículos de tração animal.

\platepair{moped-municipal-oer}{Ciclomotor de Oeiras: \plate{2-OER 20-75}.}%
          {moto-municipal-str}{Motociclo até 50~cm³ de Santarém: \plate{1-STR 00-30}.}

Com o Decreto-Lei n.º~106/2006, estas matrículas passaram a ter de ser
substituídas até 2008, ainda que algumas chapas se tenham mantido em uso até
2010.

\subsection{Motociclos e ciclomotores depois de 2006}

Os motociclos com cilindrada superior a 50~cm³ têm matrículas com a mesma
sequência das dos automóveis, mas em chapas de formato mais reduzido, sem
eurobanda e sem a barra amarela com a data da primeira matrícula.

Os motociclos com cilindrada inferior a 50~cm³ e os ciclomotores têm matrículas
de formato igual às dos motociclos com mais de 50~cm³, mas de fundo amarelo.

\platepair{moto-6060uj}{Motociclo: \plate{60-60-UJ}.}%
          {moped-08ns44}{Ciclomotor: \plate{08-NS-44}.}

\platefig[0.3\linewidth]{moto-49bx10}{Motociclo: \plate{49-BX-10}.}

\subsection{Mini carros (os “papa-reformas”)}

Têm uma matrícula com a configuração habitual, mas de fundo amarelo em vez de
branco. São ciclomotores de quatro rodas (quadriciclos): os que têm espaço para
isso podem usar uma chapa do formato dos automóveis, de fundo amarelo, em vez da
chapa dos motociclos.

\platefig{minicar-77mq39}{Mini carro: \plate{77-MQ-39}.}
